\section{Introduction}

\paragraph{Problem definition:}
Skin cancer is one of the most prevalent malignancies worldwide, with incidence 
and mortality rates rising steadily over recent decades~\cite{Wang_M}. Early and 
accurate diagnosis via dermoscopy is critical, yet the performance of data-driven 
diagnostic models is fundamentally bounded by \textbf{\emph{data scarcity in segmentation}} and \textbf{\emph{the high 
cost of pixel-level annotation}}. Public dermoscopy collections such as HAM10000~\cite{Tschandl2018HAM10000}, 
while providing a valuable starting point, offer limited pixel-level segmentation 
annotations and insufficient morphological diversity to train robust segmentation 
models. Rather than treating such collections as the final training source, 
we use them as a foundation to learn the visual statistics of dermoscopic 
images from which we can generate a far larger and more diverse annotated corpus. 

The problem defines data scarcity in dermoscopic segmentation as a bottleneck for developing accurate and generalizable models. The high cost of pixel-level annotation, which requires expert dermatologists, further exacerbates this issue. 

\paragraph{Motivation:}
Synthetic data generation offers a principled path around the scarcity of annotated
dermoscopic images: labeled samples can be produced at scale without patient privacy
concerns or the bottleneck of clinical annotation labour~\cite{du2024boosting}. This motivation
is grounded directly in our own research experience. While developing an in-house
dermoscopic lesion segmentation model based on a lightweight U-Net architecture,
we repeatedly encountered a hard ceiling: regardless of the architectural
modifications we applied, model performance stagnated due to the limited number
of available annotated samples. Recruiting expert dermatologists for additional
pixel-level mask annotation was not a feasible option at our scale.
Faced with this bottleneck, we concluded that the most practical path to breaking
through the segmentation performance barrier was not to redesign the model, but to
expand the training data itself — specifically by generating new paired image--mask
samples synthetically.

This observation led us to reframe the problem: rather than acquiring more labeled
images, we ask whether a generative model can produce realistic dermoscopic images
that \emph{conform to a given segmentation mask}. Such a model would invert the
traditional annotation pipeline — instead of experts labeling images, the model
produces images already aligned to their masks — yielding an effectively unlimited
supply of paired training data.

However, generating \emph{clinically plausible} dermoscopy images is non-trivial. 
Synthesized images must (i) respect precise lesion geometry shape, borders, and 
asymmetry that encodes diagnostic features; (ii) reproduce realistic dermoscopic 
textures such as pigment networks, streaks, and regression structures; and (iii) 
conform to the optical characteristics of real dermoscopes~\cite{Tschandl2018HAM10000}. Earlier
GAN-based approaches struggle with geometric controllability and mode
collapse~\cite{qasim2021redganattackingclassimbalance}, while recent diffusion models offer superior sample diversity
but typically lack fine-grained spatial control~\cite{SheJun_DiDGen_MICCAI2025,xu2025skindualgen}. Combining latent 
diffusion with ControlNet~\cite{zhang2023adding} provides a promising bridge: 
mask-level spatial conditioning guides lesion geometry while the pre-trained 
backbone supplies rich texture priors.
\paragraph{Contributions.}
This paper presents \textit{Mask2Derm}, a mask-conditioned latent diffusion framework for photorealistic and controllable dermoscopic image synthesis. Our main contributions are:

\begin{itemize}[leftmargin=1.5em]
    \item \textbf{Public minimal synthetic dataset.} We release a publicly available collection of synthetically generated dermoscopic images paired with binary segmentation masks, intended as a minimal ready-to-use resource for training and benchmarking segmentation models. The dataset is accessible at \href{https://www.kaggle.com/datasets/rollas/mask2derm-generated-skin-lesions}{Kaggle: Mask2Derm Generated Skin Lesions}.
    \item \textbf{Mask-conditioned ControlNet pipeline.} We fine-tune a ControlNet module on top of a frozen Stable Diffusion backbone, enabling precise geometry control from binary segmentation masks while preserving the generative diversity of the pre-trained prior.
    \item \textbf{Clinical utility validation.} [Will be added soon]
\end{itemize}
